\documentclass{article}
\usepackage{amsmath,amssymb}
\begin{document}

\section*{Higham Chapter 6, Problem 6.13}

I am formalizing Higham, \emph{Accuracy and Stability of Numerical Algorithms},
2nd ed., Chapter 6, Problem 6.13 in Lean 4 / Mathlib.

Source statement:
\[
  \text{If } H\in \mathbb{R}^{n\times n}\text{ is a Hadamard matrix, then }
  \|H\|_p=\max\{n^{1/p}, n^{1-1/p}\}.
\]
Here $\|\cdot\|_p$ is the subordinate matrix norm induced by the vector
$\ell^p$ norm, with the endpoint $p=\infty$ understood separately.

Local modeling plan:
\begin{itemize}
\item Work over finite indices \texttt{Fin n}.
\item Use a predicate roughly saying: every entry of $H$ has absolute value
  $1$, and the real columns/rows are orthogonal:
  \[
    \sum_i H_{i j} H_{i k} =
      \begin{cases} n,& j=k,\\0,& j\ne k,\end{cases}
    \quad
    \sum_j H_{i j} H_{k j} =
      \begin{cases} n,& i=k,\\0,& i\ne k.\end{cases}
  \]
\item The repository already has a local complex matrix norm API:
  \texttt{complexMatrixLpNormOfReal hn p hp A} for finite real $p\ge1$,
  \texttt{complexMatrixInfNorm A} for $p=\infty$, and
  \texttt{realRectToCMatrix H}.
\item Existing local facts include:
  \[
  \|A\|_1 = \text{max column one-norm},\quad
  \|A\|_\infty = \text{max row one-norm},
  \]
  \[
  \texttt{complexMatrixLpNormOfReal\_rieszThorin\_one\_two}
  \]
  giving the $1$--$2$ interpolation upper bound for $1\le p\le2$,
  a general endpoint-aware Riesz--Thorin theorem, finite row/column lower
  bounds, and an adjoint/transposed finite $p/q$ theorem.
\end{itemize}

Please give a rigorous proof route, with enough detail to translate into Lean,
for the full finite-real statement
\[
  \|H\|_p=\max\{n^{1/p}, n^{1-1/p}\}, \qquad 1\le p<\infty,
\]
and the endpoint $\|H\|_\infty=n$.

Please be explicit about:
\begin{enumerate}
\item the minimal Hadamard hypotheses needed;
\item the proof that $\|H\|_2=\sqrt n$ over complexified vectors from real
  orthogonality;
\item how to get the upper bound for $1\le p\le2$ and for $2\le p<\infty$;
\item how to get both lower bounds without assuming the first row/column is
  all ones;
\item any real-power identities that are likely to be needed in Lean, such as
  simplifying powers of $n$, $\sqrt n$, $1/p$, and the Holder conjugate
  exponent.
\end{enumerate}

Avoid using an unproved ``Hadamard theorem'' as a black box.  The answer is
advisory; I will translate it into checked Lean declarations and verify with
the local build.

\end{document}
